\section*{Hoofdstuk \ref{ch:b1:derivative} --- Differentiaalrekening}

\begin{solution}{exo:b1:derivative:1}
$x^x = \eu^{x\ln x}$ op $\intoo{0}{+\infty}$: \omterm{def:b1:derivative:def}{afgeleide} $(\ln x +
1)\,x^x$.

$\ln(x + \sqrt{x^2+1})$ op $\R$ (het argument is altijd $> 0$):
\omterm{def:b1:derivative:def}{afgeleide} $\frac{1}{\sqrt{x^2+1}}$ (berekend in
\cref{prop:b1:functions:invhyp} --- het is
$\operatorname{arsinh}$).

$\arctan\frac1x$ op $\R^*$: \omterm{def:b1:derivative:def}{afgeleide} $\frac{-1/x^2}{1 + 1/x^2} =
\frac{-1}{1 + x^2}$ (in overeenstemming met
\cref{prop:b1:functions:arcidentities}~(2): de functie is
$\pm\frac\pi2 - \arctan x$ op elke halfrechte).

$\sqrt{1 + \eu^{2x}}$ op $\R$: \omterm{def:b1:derivative:def}{afgeleide}
$\frac{\eu^{2x}}{\sqrt{1 + \eu^{2x}}}$.
\end{solution}

\begin{solution}{exo:b1:derivative:2}
In $0$: $\bigl|\frac{f(h) - 0}{h}\bigr| = \abs{h \sin\frac1h} \leq
\abs h \to 0$, dus $f'(0) = 0$. Voor $x \neq 0$ geven de gewone
regels $f'(x) = 2x\sin\frac1x - \cos\frac1x$. Langs $x_n =
\frac{1}{2\pi n}$: $f'(x_n) = 0 - 1 \to -1$; langs $y_n =
\frac{1}{(2n+1)\pi}$: $f'(y_n) = 0 + 1 \to 1$. Twee rijen die naar
$0$ naderen met verschillende limieten van $f'$: geen limiet
(\cref{thm:b1:continuity:seqchar}), dus $f'$ is niet \omterm{def:b1:continuity:continuous}{continu} in $0$
en $f$ is \omterm{def:b1:derivative:def}{afleidbaar} zonder $C^1$ te zijn.
\end{solution}

\begin{solution}{exo:b1:derivative:3}
$\ln(1+x) < x$ voor $x > 0$: de raaklijnongelijkheid van de
concaviteit in $0$ (strikt buiten het raakpunt, omdat $\ln$ strikt
concaaf is; of pas de middelwaardestelling toe: $\ln(1+x) =
\frac{x}{1+c}$ voor zekere $c \in \intoo{0}{x}$, en $\frac{x}{1+c} <
x$). Dezelfde middelwaarde-identiteit geeft de ondergrens:
$\frac{x}{1+c} > \frac{x}{1+x}$.

Gevolg: met $x/n$ in plaats van $x$,
\[
\frac{x/n}{1 + x/n} < \ln\Bigl(1 + \frac xn\Bigr) < \frac xn
\quad\implies\quad
\frac{x}{1 + x/n} < n \ln\Bigl(1 + \frac xn\Bigr) < x .
\]
Het linkerlid nadert tot $x$: door insluiting geldt $n\ln(1 + \frac
xn) \to x$, en wegens de \omterm{def:b1:continuity:continuous}{continuïteit} van $\exp$ is $\bigl(1 +
\frac xn\bigr)^n = \eu^{n\ln(1 + x/n)} \to \eu^x$.
\end{solution}

\begin{solution}{exo:b1:derivative:4}
Zij $x_1 < x_2 < \dots < x_k$ verschillende wortels van $P$. Op elke
$\intcc{x_i}{x_{i+1}}$ levert Rolle
(\cref{thm:b1:derivative:rolle}) een $c_i \in \intoo{x_i}{x_{i+1}}$
met $P'(c_i) = 0$: dat zijn $k - 1$ wortels van $P'$, verschillend
omdat de \omterm{def:b1:topology:open}{open} \omterm{prop:b1:reals:intervals}{intervallen} disjunct zijn --- en per constructie
afwisselend met die van $P$.

Heeft $P$ (van graad $n$) al haar wortels reëel, schrijf ze dan met
\omterm{def:b1:poly:derivative}{multipliciteiten} $m_1 + \dots + m_k = n$. Elke wortel van
\omterm{def:b1:poly:derivative}{multipliciteit} $m_i \geq 2$ is een wortel van $P'$ van
\omterm{def:b1:poly:derivative}{multipliciteit} $m_i - 1$ (\cref{prop:b1:poly:multiplicity}), goed
voor $\sum (m_i - 1) = n - k$; Rolle draagt er nog $k - 1$ bij, alle
verschillend van deze. Totaal $\geq n - 1 = \deg P'$: alle wortels
van $P'$ zijn reëel.
\end{solution}

\begin{solution}{exo:b1:derivative:5}
Leg $\varepsilon > 0$ vast en een $A$ met $\abs{f'(t) - \ell} \leq
\varepsilon$ voor $t \geq A$. Voor $x > A$ geeft de
middelwaardestelling op $\intcc{A}{x}$ een $c \in \intoo{A}{x}$ met
\[
f(x) = f(A) + f'(c)(x - A),
\qquad\text{dus}\qquad
\Bigl|\frac{f(x)}{x} - \ell\Bigr|
\leq \frac{\abs{f(A)} + \abs\ell A}{x} + \abs{f'(c) - \ell}
\cdot\frac{x - A}{x} \leq \frac{C_A}{x} + \varepsilon .
\]
Voor grote $x$ is $\frac{C_A}{x} \leq \varepsilon$: bijgevolg
$\frac{f(x)}{x} \to \ell$.

Ja: $f(x+1) - f(x) = f'(c_x)$ met $c_x \in \intoo{x}{x+1}$
(middelwaardestelling op $\intcc{x}{x+1}$), en $c_x \to +\infty$,
dus $f(x+1) - f(x) \to \ell$.
\end{solution}

\begin{solution}{exo:b1:derivative:6}
Inductie op $n$. Voor $n = 1$: Rolle. Geldt de bewering voor $n - 1$:
$f$ wordt nul in $n+1$ punten, dus wordt $f'$, volgens Rolle
toegepast op de $n$ tussenruimten, nul in $n$ verschillende punten;
de inductiehypothese toegepast op $f'$ ($n-1$ keer \omterm{def:b1:derivative:def}{afleidbaar}, $n$
nulpunten) maakt $(f')^{(n-1)} = f^{(n)}$ ergens nul.

Toepassing: wordt $P$ van graad $\leq n$ nul in $n+1$ punten, dan
wordt $P^{(n)}$, een constante gelijk aan $n!$ maal de kopcoëfficiënt,
nul: de kopcoëfficiënt is $0$, en men besluit met neerwaartse
inductie (of rechtstreeks: alle coëfficiënten zijn nul).
\end{solution}

\begin{solution}{exo:b1:derivative:7}
Volgens de middelwaardestelling op $\intcc{a}{x_0}$ en op
$\intcc{x_0}{b}$:
\[
f'(c_1) = \frac{f(x_0) - f(a)}{x_0 - a} = \frac{f(x_0)}{x_0 - a} > 0,
\qquad
f'(c_2) = \frac{f(b) - f(x_0)}{b - x_0} = \frac{-f(x_0)}{b - x_0} < 0,
\]
met $c_1 < x_0 < c_2$. De middelwaardestelling toegepast op $f'$ op
$\intcc{c_1}{c_2}$ geeft dan een $c$ met
\[
f''(c) = \frac{f'(c_2) - f'(c_1)}{c_2 - c_1} < 0 . \qedhere
\]
\end{solution}

\begin{solution}{exo:b1:derivative:8}
$f'(x) = \frac{1 - \ln x}{x^2}$: $f$ stijgt op $\intoc{0}{\eu}$,
daalt op $\intco{\eu}{+\infty}$, met maximum $f(\eu) = \frac1\eu$;
limieten $-\infty$ in $0^+$ en $0$ in $+\infty$ (vergelijking van
groeisnelheden).

Voor $\eu \leq a < b$ geeft het strikt dalen van $f$ daar dat
$\frac{\ln a}{a} > \frac{\ln b}{b}$, dat wil zeggen $b \ln a > a \ln
b$, dus $a^b > b^a$.

Met $a = \eu < b = \pi$: $\eu^\pi > \pi^\eu$.

Kleine paren: $(2, 3)$: $2^3 = 8 < 9 = 3^2$ --- omgekeerd! De reden:
$2 < \eu$, en op $\intoo{0}{\eu}$ is de functie $f$ \emph{stijgend},
zodat de vergelijking omklapt wanneer beide getallen onder $\eu$
liggen, en onvoorspelbaar is wanneer ze $\eu$ omspannen ($f(2) =
f(4)$ verklaart het gelijkspel $2^4 = 4^2 = 16$).
\end{solution}

\begin{solution}{exo:b1:derivative:9}
\emph{Jensen voor $\ln$, met inductie op $n$.} Bewering: voor
positieve $x_i$ en gewichten $\lambda_i > 0$ met $\sum \lambda_i =
1$ geldt $\ln\bigl(\sum \lambda_i x_i\bigr) \geq \sum \lambda_i \ln
x_i$. Voor $n = 2$ is dit de concaviteit. Stap: met $\Lambda =
\lambda_1 + \dots + \lambda_{n-1} = 1 - \lambda_n$ en $y = \sum_{i<n}
\frac{\lambda_i}{\Lambda} x_i$ is
\[
\ln\Bigl(\sum_{i \leq n} \lambda_i x_i\Bigr)
= \ln\bigl(\Lambda y + \lambda_n x_n\bigr)
\geq \Lambda \ln y + \lambda_n \ln x_n
\geq \Lambda \sum_{i<n} \frac{\lambda_i}{\Lambda}\ln x_i
+ \lambda_n \ln x_n,
\]
met gebruik van de concaviteit ($n = 2$) en daarna de
inductiehypothese.

Met $\lambda_i = \frac 1n$ en $x_i = a_i$: $\ln\frac{\sum a_i}{n}
\geq \frac 1n \sum \ln a_i = \ln\sqrt[n]{a_1\cdots a_n}$; neem de
exponentiële. Gelijkheid: $\ln$ is \emph{strikt} concaaf ($\ln'' <
0$), zodat gelijkheid bij elke stap afdwingt dat de gemiddelde
punten samenvallen --- dat wil zeggen dat alle $a_i$ gelijk zijn; en
zijn ze alle gelijk, dan is de gelijkheid duidelijk.
\end{solution}

\begin{solution}{exo:b1:derivative:10}
Zij $g(x) = f(x) - vx$: \omterm{def:b1:derivative:def}{afleidbaar}, met $g'(a) = f'(a) - v < 0$ en
$g'(b) = f'(b) - v > 0$. Volgens de extremumstelling bereikt $g$ haar
minimum op $\intcc{a}{b}$ in zeker punt $c$. Dat is niet $a$: omdat
$g'(a) < 0$, geldt $g < g(a)$ net rechts van $a$. Het is niet $b$:
omdat $g'(b) > 0$, geldt $g < g(b)$ net links van $b$. Dus is $c$
\omterm{def:b1:topology:closure}{inwendig}, en geeft \cref{prop:b1:derivative:fermat} dat $g'(c) = 0$,
dat wil zeggen $f'(c) = v$.
\end{solution}

\begin{solution}{exo:b1:derivative:11}
\emph{Eenduidigheid:} twee vaste punten $\ell \neq \ell'$ zouden
$\abs{\ell - \ell'} = \abs{f(\ell) - f(\ell')} \leq k\abs{\ell -
\ell'} < \abs{\ell - \ell'}$ geven, absurd.

\emph{Bestaan:} $g(x) = f(x) - x$ voldoet, wegens de
middelwaarde-ongelijkheid, aan $f(x) \leq f(0) + k\abs x$; dus voor
$x \geq \frac{\abs{f(0)}}{1 - k}$ is $g(x) \leq f(0) + kx - x \leq
0$, en symmetrisch is $g(-x) \geq 0$ voor grote $x$. De
tussenwaardestelling geeft een nulpunt $\ell$ van $g$: een vast punt.

\emph{Convergentie:} opnieuw de middelwaarde-ongelijkheid:
\[
\abs{u_{n+1} - \ell} = \abs{f(u_n) - f(\ell)} \leq k\abs{u_n - \ell},
\]
dus met inductie $\abs{u_n - \ell} \leq k^n \abs{u_0 - \ell} \to 0$.
\end{solution}

\begin{solution}{exo:b1:derivative:12}
\begin{enumerate}
  \item Als $g(b) = g(a)$, zou Rolle een \omterm{def:b1:topology:closure}{inwendig} nulpunt van
        $g'$ geven: uitgesloten. Stel $\lambda = \frac{f(b) -
        f(a)}{g(b) - g(a)}$ en $h = f - \lambda g$: $h$ is \omterm{def:b1:continuity:continuous}{continu}
        op $\intcc{a}{b}$, \omterm{def:b1:derivative:def}{afleidbaar} binnenin, en $h(b) - h(a) =
        f(b) - f(a) - \lambda(g(b) - g(a)) = 0$. Rolle levert een
        $c$ met $h'(c) = 0$, dat wil zeggen $f'(c) =
        \lambda\,g'(c)$; deel door $g'(c) \neq 0$.
  \item Voor $x > a$ dicht bij $a$ geeft deel (1) op
        $\intcc{a}{x}$ (waar $g' \neq 0$) dat $g(x) \neq 0$ en een
        $c_x \in \intoo{a}{x}$ met
        \[
        \frac{f(x)}{g(x)} = \frac{f(x) - f(a)}{g(x) - g(a)}
        = \frac{f'(c_x)}{g'(c_x)} .
        \]
        Als $x \to a^+$, dan $c_x \to a^+$ (insluiting), dus
        nadert het rechterlid tot $\ell$: $\frac{f}{g} \to \ell$.
  \item $\frac{f(x)}{g(x)} = x\sin\frac1x \to 0$, terwijl
        $\frac{f'(x)}{g'(x)} = 2x\sin\frac1x - \cos\frac1x$ geen
        limiet heeft in $0$ (\cref{exo:b1:derivative:2}): de regel
        van l'Hôpital draagt informatie alleen over van
        $\frac{f'}{g'}$ naar $\frac fg$, nooit terug.
\end{enumerate}
\end{solution}

\begin{solution}{pb:b1:derivative:1}
\textbf{1.} $\bigl|\frac ab - \frac pq\bigr| = \frac{\abs{aq -
bp}}{bq}$, en $aq - bp$ is een niet-nul geheel getal wanneer de
breuken verschillen: de afstand is $\geq \frac{1}{bq}$. Dus komt
geen enkel rationaal getal behalve $x$ zelf in het geperforeerde
\omterm{prop:b1:reals:intervals}{interval} met straal $\frac{1}{bq}$ rond $x = \frac ab$.

\textbf{2.} Is $\bigl|\sqrt2 - \frac pq\bigr| \geq 1 \geq
\frac{1}{4q^2}$, dan zijn we klaar. Anders is $\frac pq \in
\intoo{\sqrt2 - 1}{\sqrt2 + 1}$, dus $0 < \sqrt2 + \frac pq <
2\sqrt2 + 1 < 4$. Omdat $\sqrt 2 \notin \Q$, is $p^2 - 2q^2$ een
niet-nul geheel getal, en
\[
\Bigl|\sqrt2 - \frac pq\Bigr|
= \frac{\abs{2q^2 - p^2}}{q^2\,\bigl(\sqrt2 + \frac pq\bigr)}
\geq \frac{1}{4q^2} .
\]

\textbf{3.} $(p + 2q)^2 - 2(p + q)^2 = -(p^2 - 2q^2)$: de waarde
$\pm1$ plant zich voort. Vanuit $(1,1)$:
\[
(3,2),\ (7,5),\ (17,12),\ (41,29),\ (99,70),
\]
waarbij $p^2 - 2q^2$ afwisselt tussen $-1$ en $+1$. Voor deze
paren is $\frac pq \geq 1$, dus $\sqrt2 + \frac pq > 2$ en
\[
\Bigl|\sqrt2 - \frac pq\Bigr| =
\frac{1}{q^2(\sqrt2 + p/q)} < \frac{1}{2q^2} ,
\]
met $q \to \infty$: oneindig veel benaderingen van orde $2$. Met
vraag 2 is de exponent $2$ exact voor $\sqrt 2$.

\textbf{4.} De $N + 1$ getallen $kx - \lfloor kx\rfloor$ ($0
\leq k \leq N$) liggen in de $N$ hokjes $\intco{\frac
jN}{\frac{j+1}{N}}$: twee ervan delen een hokje
(\cref{cor:b1:counting:pigeonhole}), zeg voor $i < j$. Met $q =
j - i \leq N$ en $p = \lfloor jx\rfloor - \lfloor ix\rfloor$
geldt $\abs{qx - p} < \frac1N$, dus $\bigl|x - \frac pq\bigr| <
\frac{1}{Nq} \leq \frac{1}{q^2}$. Laat $N \to \infty$: omdat $x$
irrationaal is, heeft elke vaste breuk een positieve afstand tot
$x$, terwijl $\frac{1}{Nq} \leq \frac 1N \to 0$ afdwingt dat er
nieuwe breuken opduiken: oneindig veel verschillende $\frac pq$
met $\bigl|x - \frac pq\bigr| < \frac{1}{q^2}$.

\textbf{5.} Vertrek van een willekeurige niet-nulle gehele $P_0$
met $P_0(x) = 0$. Heeft $P_0$ een rationale wortel $\frac ab$, dan
schrijft de factorstelling (\cref{thm:b1:poly:factor}) $P_0 =
\bigl(X - \frac ab\bigr)Q$ met $Q \in \Q[X]$; omdat $x \neq \frac
ab$ ($x$ is irrationaal), is $Q(x) = 0$, en het wegwerken van de
noemers geeft een niet-nulle \emph{gehele} \omterm{def:b1:poly:def}{veelterm} van kleinere
graad die in $x$ nul wordt. De graad daalt bij elke stap, dus het
proces stopt: we bereiken $P \in \Z[X]$, $P(x) = 0$, zonder
rationale wortel, van zekere graad $d$. Was $d \leq 1$, dan zou $P
= uX + v$ het getal $x = -\frac vu$ rationaal maken: dus $d \geq
2$.

\textbf{6.} $q^d\,P\bigl(\frac pq\bigr) = a_d p^d + a_{d-1}
p^{d-1} q + \dots + a_0 q^d$ is een geheel getal, en het is niet
nul omdat $P$ geen rationale wortel heeft:
$\bigl|P\bigl(\frac pq\bigr)\bigr| \geq q^{-d}$.

\textbf{7.} Merk op dat $M > 0$: $P'$ is een niet-nulle \omterm{def:b1:poly:def}{veelterm}
($d \geq 2$), dus zij kan niet identiek nul zijn op
$\intcc{x-1}{x+1}$. Is $\bigl|x - \frac pq\bigr| > 1$, dan
overtreft dit $\frac{C}{q^d}$ triviaal. Anders is $\frac pq \in
\intcc{x-1}{x+1}$ en geeft de middelwaardestelling
(\cref{thm:b1:derivative:mvt}) een $c$ tussen $x$ en $\frac pq$
met
\[
\Bigl|P\Bigl(\frac pq\Bigr)\Bigr|
= \Bigl|P\Bigl(\frac pq\Bigr) - P(x)\Bigr|
= \abs{P'(c)}\,\Bigl|x - \frac pq\Bigr|
\leq M\,\Bigl|x - \frac pq\Bigr| ,
\]
dus met vraag 6: $\bigl|x - \frac pq\bigr| \geq \frac{1}{Mq^d}
\geq \frac{C}{q^d}$.

\textbf{8.} Stel dat $x = \frac ab$ een Liouville-getal is. Kies
$n$ met $2^{n-1} > b$ en het bijbehorende $\frac pq$, $q \geq 2$:
\[
0 < \Bigl|x - \frac pq\Bigr| < \frac{1}{q^n}
= \frac{1}{q^{n-1}\,q} \leq \frac{1}{2^{n-1} q} <
\frac{1}{bq} ,
\]
in tegenspraak met vraag 1. Liouville-getallen zijn dus
irrationaal.

\textbf{9.} Was een Liouville-getal $x$ algebraïsch, dan is het
irrationaal (vraag 8), zodat de vragen 5--7 een $d \geq 2$ en een
$C > 0$ leveren met steeds $\bigl|x - \frac pq\bigr| \geq
\frac{C}{q^d}$. Voor elke $n$ geeft de Liouville-benadering dat
$\frac{C}{q^d} < q^{-n}$, dat wil zeggen $C < q^{d-n} \leq
2^{d-n}$ (want $q \geq 2$). Voor grote $n$ is $2^{d-n} < C$:
tegenspraak. Liouville-getallen zijn \omterm{pb:b1:logic:1}{transcendent}.

\textbf{10.} Enen op de posities $1, 2, 6, 24$; alle andere
cijfers onder de eerste $25$ zijn nul:
\[
L = 0.1100010000\,0000000000\,00010\dots
\]

\textbf{11.} Voor $m > k$ zijn de posities $n!$ met $n > k$
verschillende gehele getallen $\geq (k+1)!$, dus geeft de eindige
meetkundige som
\[
t_m - t_k = \sum_{n=k+1}^{m} 10^{-n!}
\leq \sum_{j = (k+1)!}^{m!} 10^{-j}
< 10^{-(k+1)!}\,\frac{1}{1 - \frac1{10}}
= \frac{10}{9}\,10^{-(k+1)!} ;
\]
neem het \omterm{def:b1:reals:bounds}{supremum} over $m$: $L - t_k \leq
\frac{10}{9}10^{-(k+1)!} < 2\cdot10^{-(k+1)!}$. Ondergrens: $L
\geq t_{k+1} = t_k + 10^{-(k+1)!}$.

\textbf{12.} $p_k = 10^{k!}\,t_k \in \N$, $q_k = 10^{k!}$, en
$(k+1)! = (k+1)\,k!$ geeft $10^{-(k+1)!} = q_k^{-(k+1)}$: vraag
11 luidt
\[
0 < L - \frac{p_k}{q_k} < \frac{2}{q_k^{\,k+1}} .
\]
Zij $n$ gegeven: voor $k \geq n$ is $2\,q_k^{-(k+1)} \leq
q_k^{-n}$ (immers $q_k^{\,k+1-n} \geq q_k \geq 10 > 2$), en $q_k
\geq 2$: aan de definitie van vraag 8 is voldaan. $L$ is een
Liouville-getal.

\textbf{13.} Volgens vraag 9 is $L$ \omterm{pb:b1:logic:1}{transcendent} --- het eerste
getal in de geschiedenis waarvan de transcendentie werd bewezen
(Liouville, 1844). Controle op de cijfers: de reeks heeft
oneindig veel enen met opeenvolgende gaten $(k+1)! - k! = k\cdot
k! \to \infty$, dus is zij niet uiteindelijk periodiek, en is $L
\notin \Q$ volgens het periodiciteitscriterium van
\cref{pb:b1:reals:1} (vraag 18) --- consistent.

\textbf{14.} Met cijfers $d_n \in \intint{1}{9}$ op de
faculteitsposities: de staartgrens van vraag 11 schaalt met
hoogstens $9$: $0 < L' - t'_k \leq
9\cdot\frac{10}{9}\,10^{-(k+1)!} = 10\,q_k^{-(k+1)}$ (positief
omdat het cijfer op positie $(k+1)!$ niet nul is). Voor $k \geq
n$ is $10\,q_k^{-(k+1)} \leq q_k^{-n}$, want $q_k^{\,k+1-n} \geq
10$: opnieuw een Liouville-getal, dus \omterm{pb:b1:logic:1}{transcendent}. Deze waarden
zijn paarsgewijs verschillend voor verschillende cijferkeuzen (de
reeksen zijn eigenlijk --- nullen zijn er in overvloed --- en
eigenlijke reeksen bepalen hun waarde, \cref{pb:b1:reals:1},
vraag 10). Is een lijst $k \mapsto x_k$ van zulke getallen
gegeven, kies dan het $k$-de faculteitscijfer in $\intint{1}{9}$
verschillend van dat van $x_k$: een getal van dezelfde vorm dat
in de lijst ontbreekt. Overaftelbaar veel expliciete
\omterm{pb:b1:logic:1}{transcendente} getallen.

\textbf{15.} Eerst een lemma: \emph{geldt $\bigl|x - \frac
pq\bigr| \geq \frac{C}{q^s}$ voor alle $\frac pq \neq x$, dan is
$x$ tot geen enkele orde $\mu > s$ benaderbaar.} Immers,
oneindig veel $\frac pq \neq x$ met $\bigl|x - \frac pq\bigr| <
\frac{c}{q^\mu}$ zouden $\frac{C}{q^s} < \frac{c}{q^\mu}$
afdwingen, dat wil zeggen $q^{\mu - s} < \frac cC$: de $q$ zijn
begrensd, en begrensd veel breuken liggen binnen afstand $1$ van
$x$ --- eindig veel kandidaten, geen oneindig veel. Nu de
hiërarchie: rationale getallen zijn benaderbaar tot orde $1$
($\frac pq$ met $p = \lfloor qx\rfloor + 1$ geeft een fout $\leq
\frac1q < \frac2q$) en tot geen enkele orde $\mu > 1$ (vraag 1
geeft de hypothese van het lemma met $s = 1$, $C = \frac1b$);
$\sqrt2$: orde $2$ (vraag 3) en niet meer (vraag 2 en het lemma);
elk irrationaal getal: minstens $2$ (vraag 4); een \omterm{pb:b1:logic:1}{algebraïsch
getal} van graad $d$: hoogstens $d$ (vraag 7 en het lemma);
Liouville-getallen: elke orde (de formule van vraag 12, met $c =
2$).

\textbf{16.} Met $r = \frac ab$: $\frac{p_k}{q_k} + \frac ab =
\frac{b p_k + a q_k}{b q_k} =: \frac{P_k}{Q_k}$, $Q_k = b q_k
\geq 2$, en
\[
\Bigl|(L + r) - \frac{P_k}{Q_k}\Bigr| = L - \frac{p_k}{q_k}
< 2\,q_k^{-(k+1)} = 2\,b^{\,k+1} Q_k^{-(k+1)} .
\]
Zij $n$ gegeven: voor grote $k$ is $Q_k^{\,k+1-n} \geq Q_k =
b\,10^{k!} \geq 2\,b^{\,k+1}$ (de faculteit verplettert de
macht), zodat de fout $< Q_k^{-n}$ is: $L + r$ is een
Liouville-getal. Omdat $\Q$ \omterm{def:b1:topology:dense}{dicht} ligt en elke $L + r$
\omterm{pb:b1:logic:1}{transcendent} is, liggen de \omterm{pb:b1:logic:1}{transcendente} getallen \omterm{def:b1:topology:dense}{dicht} in $\R$.

\textbf{17.} Voor $h \geq 1$ zijn er eindig veel $P \in \Z[X]$
met $\deg P + \sum_i \abs{a_i} \leq h$ (graad $\leq h$ en elke
coëfficiënt in $\intint{-h}{h}$: hoogstens $(2h+1)^{h+1}$). Elke
niet-nulle gehele \omterm{def:b1:poly:def}{veelterm} heeft zo'n hoogte, en heeft hoogstens
$\deg P$ reële wortels: de algebraïsche getallen vormen een
aftelbare vereniging (over $h$) van \omterm{def:b1:counting:card}{eindige verzamelingen}, en
kunnen dus als één rij worden opgesomd. Konden ook de
\omterm{pb:b1:logic:1}{transcendente} getallen worden opgesomd, dan zou het verweven van
de twee lijsten $\R$ opsommen, in tegenspraak met
\cref{pb:b1:reals:1} (vraag 22). De \omterm{pb:b1:logic:1}{transcendente} getallen vormen
dus een overaftelbare \omterm{def:b1:logic:sets}{verzameling}. Vergelijking: Cantor bewijst
dat de \emph{meeste} reële getallen \omterm{pb:b1:logic:1}{transcendent} zijn maar toont
er geen enkel; Liouville toont er één, met effectieve constanten
(deel V) --- bestaan door overvloed tegenover bestaan door
constructie.

\textbf{18.} Uit vraag 2 geldt de hypothese van het lemma met $s
= 2$, $C = \frac14$. Was $\sqrt2$ een Liouville-getal, dan geldt
voor $n = 3$: $\frac{1}{4q^2} < q^{-3}$ dwingt $q < 4$ af, dus $q
\in \{2, 3\}$; slechts eindig veel $\frac pq$ met deze $q$ liggen
binnen afstand $1$ van $\sqrt2$, elk op een positieve afstand
$\geq \varepsilon_0$ ($\sqrt2$ is irrationaal); door $n$ te
kiezen met $2^{-n} < \varepsilon_0$ blijft er helemaal geen
toelaatbare $\frac pq$ over: tegenspraak. Hetzelfde argument met
$\frac{C}{q^d}$ toont aan dat geen enkel \omterm{pb:b1:logic:1}{algebraïsch getal} een
Liouville-getal is --- vraag 9 in effectieve kleren.

\textbf{19.} $99^2 - 2\cdot70^2 = 9801 - 9800 = 1$. Bijgevolg
\[
\sqrt2 - \frac{99}{70} =
\frac{2 - (99/70)^2}{\sqrt2 + 99/70}
= \frac{-1}{4900\,\bigl(\sqrt2 + \tfrac{99}{70}\bigr)} ,
\qquad
\Bigl|\sqrt2 - \frac{99}{70}\Bigr|
= \frac{1}{4900 \times 2.8284\dots} \approx 7.2\cdot10^{-5} :
\]
$\frac{99}{70} = 1.414285\dots$ tegenover $\sqrt2 =
1.414213\dots$ --- vijf juiste cijfers.

\textbf{20.} Test op rationale wortels voor $P = X^3 - 2$:
kandidaten $\pm1, \pm2, \pm\frac12$, geen ervan is een wortel.
Dus $d = 3$ en deel II is van toepassing op $x = 2^{1/3} =
1.2599\dots$ Op $\intcc{x - 1}{x + 1} \subseteq \intcc{0.25}{2.26}$
is $\abs{P'(t)} = 3t^2 \leq 3\,(1 + 2^{1/3})^2 < 3\times(2.26)^2 =
15.32 < 16$, dus $M < 16$ en $C \geq \frac{1}{16}$:
\[
\Bigl|2^{1/3} - \frac pq\Bigr| \geq \frac{1}{16\,q^3}
\qquad\text{voor alle rationale getallen.}
\]

\textbf{21.} Is $\bigl|2^{1/3} - \frac pq\bigr| < 10^{-6}$, dan
is $\frac{1}{16 q^3} < 10^{-6}$, dat wil zeggen $q^3 >
\frac{10^6}{16} = 62\,500$; omdat $39^3 = 59\,319 < 62\,500 \leq
64\,000 = 40^3$, is $q \geq 40$.

\textbf{22.} Laat deel III lopen in grondtal $2$: $B = \sup_k
\sum_{n\leq k} 2^{-n!}$, $q_k = 2^{k!}$, en de meetkundige staart
(reden $\frac12$) geeft $0 < B - \frac{p_k}{q_k} \leq
2\cdot2^{-(k+1)!} = 2\,q_k^{-(k+1)} \leq q_k^{-n}$ voor $k \geq
n$. Dus is $B$ een Liouville-getal, en bijgevolg \omterm{pb:b1:logic:1}{transcendent}:
niets in het argument is decimaal.

\textbf{23.} Met enen op de posities $3^k$: $q_k = 10^{3^k}$ en
de staartgrens geeft $0 < x^\dagger - \frac{p_k}{q_k} <
2\cdot10^{-3^{k+1}} = 2\,q_k^{-3}$ (want $3^{k+1} = 3\cdot3^k$):
oneindig veel benaderingen van orde $3$. Volgens het lemma van
vraag 15 sluit orde $3 > 1$ rationaliteit uit, en sluit orde $3 >
2$ uit dat het een kwadratisch irrationaal getal is (waarvan de
ongelijkheid van Liouville $s = d = 2$ heeft). Maar een
\omterm{pb:b1:logic:1}{algebraïsch getal} van graad $\geq 3$ wordt slechts op orde $d
\geq 3$ afgestoten: de methode van Liouville kan $x^\dagger$ niet
van de derdegraadsgetallen scheiden. De kloof wordt gedicht door
de stelling van Roth --- elk algebraïsch irrationaal getal heeft
benaderingsorde precies $2$ --- een resultaat uit de twintigste
eeuw, ver voorbij dit volume; aangenomen dat het geldt, is ook
$x^\dagger$ \omterm{pb:b1:logic:1}{transcendent}.

\textbf{24.} Er zijn hoogstens $(2H+1)^{d+1}$ tupels $(a_0,
\dots, a_d)$ met ingangen in $\intint{-H}{H}$, en elke niet-nulle
\omterm{def:b1:poly:def}{veelterm} daaronder heeft hoogstens $d$ reële wortels: er ontstaan
hoogstens $d\,(2H+1)^{d+1}$ algebraïsche getallen --- de
eindigheid die vraag 17 in staat stelde ze alle op te sommen.

\textbf{25.} (i) Het enige analytische ingrediënt is de
middelwaardestelling, in vraag 7, die het nul worden $P(x) = 0$
omzet in de lipschitzafstoting $\abs{P(p/q)} \leq M\abs{x -
p/q}$. (ii) De spanning: de geheeltalligheid duwt $\abs{P(p/q)}$
omhoog tot $q^{-d}$, de gladheid trekt het omlaag tot $M\abs{x -
p/q}$ --- een rationaal getal te dicht bij $x$ zou tussen beide
worden verpletterd. (iii) Liouville construeert één \omterm{pb:b1:logic:1}{transcendent
getal} met effectieve constanten; Cantor toont aan dat bijna alle
reële getallen \omterm{pb:b1:logic:1}{transcendent} zijn zonder er één te noemen:
constructie tegenover \omterm{def:b1:counting:card}{kardinaliteit}. (iv) De weekendopgave van
\cref{ch:b1:integration} bewijst de irrationaliteit van $\pi$ met
dezelfde wurggreep --- een integraal die een positief geheel
getal zou zijn maar gevangen zit in $\intoo{0}{1}$ --- met
integratie in plaats van differentiatie als analytische helft.
\end{solution}
